% =========================================================================
%  Figura: geometría de un sistema de conductores. Punto fuente r'_j sobre el
%  conductor j, punto r_i sobre el conductor i y punto de observación r;
%  R_j = r - r'_j y R_ij = r_i - r'_j.
% =========================================================================
\begin{figure}[t]
\centering
\begin{tikzpicture}[font=\small, x=1cm, y=1cm]
  % ejes
  \draw[eje] (0,0) -- (0,1.6) node[above, font=\footnotesize] {$z$};
  \draw[eje] (0,0) -- (1.3,0) node[right, font=\footnotesize] {$y$};
  \draw[eje] (0,0) -- (-0.75,-0.6) node[below left=-2pt, font=\footnotesize] {$x$};
  \node[below=1pt, font=\footnotesize] at (0,0) {$O$};
  % conductor j
  \fill[top color=azulvelo!80!white, bottom color=metal, draw=azulprof, line width=0.6pt]
    (1.0,2.0) .. controls (0.8,2.9) and (1.9,3.5) .. (2.5,2.9)
    .. controls (3.1,2.4) and (2.6,1.4) .. (1.9,1.35) .. controls (1.4,1.3) and (1.1,1.5) .. cycle;
  \node[font=\footnotesize, text=azulprof] at (1.1,3.35) {$S_j'$};
  \node[font=\footnotesize] at (2.25,2.55) {$Q_j$};
  \node[font=\footnotesize, text=cargapos] at (3.55,3.35) {$\sigma_j(\rp_j)$};
  % conductor i
  \fill[top color=azulvelo!80!white, bottom color=metal, draw=azulprof, line width=0.6pt]
    (6.1,0.2) .. controls (5.7,0.9) and (6.6,1.4) .. (7.3,1.0)
    .. controls (7.9,0.6) and (7.6,-0.7) .. (7.0,-0.9) .. controls (6.5,-1.05) and (6.4,-0.4) .. cycle;
  \node[font=\footnotesize, text=azulprof] at (7.95,0.95) {$S_i'$};
  \node[font=\footnotesize] at (7.0,-0.35) {$Q_i$};
  % tercer conductor
  \fill[top color=azulvelo!80!white, bottom color=metal, draw=azulprof, line width=0.6pt]
    (7.2,3.0) .. controls (7.0,3.6) and (7.8,3.9) .. (8.2,3.5) .. controls (8.6,3.1) and (8.1,2.5) .. (7.7,2.6) .. controls (7.4,2.65) .. cycle;
  \node[font=\footnotesize] at (7.75,3.15) {$Q_N$};
  % vectores
  \coordinate (rpj) at (1.95,1.36);
  \coordinate (ri) at (6.25,0.45);
  \coordinate (r) at (5.3,2.4);
  \draw[vect, azulacero] (0,0) -- (rpj) node[midway, left=1pt, font=\footnotesize] {$\rp_j$};
  \draw[vect, azulacero] (0,0) -- (ri) node[pos=0.6, below=1pt, font=\footnotesize] {$\vc{r}_i$};
  \draw[vect, azulacero] (0,0) -- (r) node[pos=0.62, above=1pt, font=\footnotesize] {$\vc{r}$};
  \draw[vectrojo] (rpj) -- (r) node[pos=0.62, above=1pt, font=\footnotesize] {$\vc{R}_j$};
  \draw[vectrojo] (rpj) -- (ri) node[pos=0.72, below left=0pt, font=\footnotesize] {$\vc{R}_{ij}$};
  \fill[azulprof] (rpj) circle (1.3pt) (ri) circle (1.3pt) (r) circle (1.3pt);
\end{tikzpicture}
\caption{Sistema de conductores finitos a distancias finitas. El conductor $j$ tiene carga $Q_j$ repartida con densidad $\sigma_j$ sobre su superficie
$S_j'$; $\vc{R}_j=\vc{r}-\rp_j$ va de un punto fuente al punto de observación, y $\vc{R}_{ij}=\vc{r}_i-\rp_j$, de un punto fuente de $j$ a un punto del
conductor $i$.}
\label{fig:sistema-geom}
\end{figure}
